\section{Setup: A Scale Family of Teacher Targets}
\label{sec:setup}

Let $\mathcal{D}=\{x_i\}_{i=1}^{N}$ be a training corpus, $T$ a frozen teacher encoder
producing $t_i = T(x_i)\in\R^{d_T}$, and $S_\theta$ a compact student producing
$s_i = S_\theta(x_i)\in\R^{d_S}$. Teacher embeddings are computed once, cached, and the
teacher discarded; only cosine similarities of teacher vectors are ever used, so nothing is
assumed about $d_T$, the teacher's tokenizer, or its internal states.

\paragraph{Teacher affinity.}
For a temperature $\tau>0$ define the affinity kernel
\begin{equation}
    K^{(\tau)}_{ij} = \exp\!\big(\cos(t_i,t_j)/\tau\big)\,\1_{i\neq j},
    \label{eq:kernel}
\end{equation}
which is dense: every ordered pair receives positive affinity. A neighbourhood graph is
obtained by masking $K$ to mutual $k$-nearest neighbours,
\begin{equation}
    M_{ij} = \1\!\left[\,j\in\mathcal{N}_k(i)\ \wedge\ i\in\mathcal{N}_k(j)\,\right],
    \qquad
    W = M \odot K^{(\tau_g)},
    \label{eq:graph}
\end{equation}
where $\mathcal{N}_k(i)$ is the teacher's top-$k$ neighbour set of $i$. Mutuality is the
standard defence against hubness \citep{radovanovic2010hubs}; Section~\ref{sec:support}
shows it has a cost. With $D=\mathrm{diag}(W\mathbf{1})$ we write $P = D^{-1}W$ for the
row-stochastic transition matrix and $Q = \tfrac{1}{2}(I + P)$ for its lazy variant.

\paragraph{The scale family.}
Sparsification and diffusion are two separate operations on the same kernel, and
indexing them jointly yields a single family. Let
$\Rset = \{0\}\cup\mathcal{R}$ with $\mathcal{R}\subset\mathbb{N}_{>0}$, and define
\begin{equation}
    P_r =
    \begin{cases}
        \rownorm\!\left(K^{(\tau_0)}\right), & r = 0 \quad\text{(no sparsification, no diffusion)},\\[2pt]
        Q^{\,r}, & r \in \mathcal{R} \quad\text{($r$-step lazy walk on the sparsified kernel)}.
    \end{cases}
    \label{eq:family}
\end{equation}
The index $r$ therefore reads as \emph{how much graph structure has been imposed}. At $r=0$
the teacher is asked directly: row $i$ of $P_0$ is its own similarity distribution over the
whole corpus. At $r\geq 1$ the kernel is first pruned to a manifold approximation, then mass
is propagated along it, so larger $r$ describes broader semantic regions. The two extremes
are not different kinds of supervision but the endpoints of one construction, which is why
\method needs no second loss term and no balancing weight.

Laziness is what makes $\mathcal{R}$ graded rather than a set of near-duplicates. Since
$\spec(P)\subseteq[-1,1]$, the map $\lambda\mapsto\frac{1}{2}(1+\lambda)$ sends $\spec(Q)$
into $[0,1]$, removing sign oscillation and roughly doubling the mixing time. Scales stay
distinguishable only while $|\lambda_2(Q)|^{r}$ has not decayed, which for a plain walk on a
mutual $k$NN graph fails within about two steps.

\paragraph{Candidate sets and restricted targets.}
Scoring an anchor against all $N-1$ columns is infeasible, so each anchor $i$ is scored
against a candidate set $C_i\subset[N]\setminus\{i\}$ with $|C_i| = C \ll N$. The target
at scale $r$ is the corresponding row of $P_r$ restricted and renormalised,
\begin{equation}
    p_{i,r}(j) = \frac{(P_r)_{ij}}{\sum_{j'\in C_i}(P_r)_{ij'}},
    \qquad j \in C_i .
    \label{eq:target}
\end{equation}
For $r\geq1$ we additionally retain only a sparse pool $\mathcal{P}_i$ of size $P$ per
anchor, whose truncation residual
$\varepsilon_{i,r} = 1-\sum_{j\in\mathcal{P}_i}(P_r)_{ij}$ is reported rather than assumed
negligible. Note $\supp(p_{i,0}) = C_i$ while $\supp(p_{i,r})\subseteq \mathcal{P}_i\cap C_i$
for $r\geq1$; this asymmetry is the subject of Section~\ref{sec:support}.

\paragraph{Student readout.}
Writing $v_i\in\R^{C}$ for the student's cosine profile, $v_i(j) = \cos(s_i,s_j)$, a
readout at temperature $\tau$ is
\begin{equation}
    q_{i}^{\tau}(j) = \frac{\exp\big(v_i(j)/\tau\big)}{\sum_{j'\in C_i}\exp\big(v_i(j')/\tau\big)} .
    \label{eq:readout}
\end{equation}
Every objective below is a weighted sum of $\KL(p_{i,r}\Vert q_i^{\tau_r})$ over $r\in\Rset$
with $\omega_r\geq0$, $\sum_r\omega_r=1$. The only design question is how many distinct
temperatures the readout uses, and Section~\ref{sec:analysis} shows that this single question
decides whether the objective is multi-resolution at all, whether it constrains similarity
levels, and whether it agrees with the teacher off the graph.
